% ============================================================
%  Lernzettel Template (my_dhbw Stil, analysis.tex)
%  Eine Spalte, mit TOC, accent-Theme.
%  Renderer: pdflatex (zweimal)
% ============================================================
\documentclass[11pt, a4paper]{article}

\usepackage[ngerman]{babel}
\usepackage{fontspec}
\setmainfont[
  Path=/usr/share/fonts/TTF/,
  Extension=.ttf,
  BoldFont=DejaVuSans-Bold,
  ItalicFont=DejaVuSans-Oblique,
  BoldItalicFont=DejaVuSans-BoldOblique
]{DejaVuSans}
\setsansfont[
  Path=/usr/share/fonts/TTF/,
  Extension=.ttf,
  BoldFont=DejaVuSans-Bold,
  ItalicFont=DejaVuSans-Oblique,
  BoldItalicFont=DejaVuSans-BoldOblique
]{DejaVuSans}
\setmonofont[Path=/usr/share/fonts/TTF/,Extension=.ttf]{DejaVuSansMono}
\usepackage[top=2cm, bottom=2cm, left=2cm, right=2cm, footskip=1cm]{geometry}
\usepackage{amsmath, amssymb}
\usepackage{xcolor}
\usepackage{titlesec}
\usepackage{titletoc}
\usepackage{enumitem}
\usepackage{fancyhdr}
\usepackage{lastpage}
\usepackage{microtype}
\usepackage{parskip}

\usepackage{booktabs}
\usepackage{array}
\usepackage{longtable}
\usepackage{tabularx}
\usepackage{ltablex}
\keepXColumns
\usepackage{colortbl}
\usepackage[most,breakable]{tcolorbox}
\usepackage{listings}
\usepackage[normalem]{ulem}
\usepackage{pdflscape}
\usepackage[colorlinks=true, linkcolor=accent, urlcolor=accent]{hyperref}

\renewcommand{\familydefault}{\sfdefault}

% ── Theme (my_dhbw accent) ────────────────────────────────────
\definecolor{accent}{RGB}{0, 60, 113}
\definecolor{primaryblue}{RGB}{0, 60, 113}
\definecolor{codebg}{RGB}{245, 245, 245}
\definecolor{figurebg}{RGB}{232, 241, 255}
\definecolor{tablehintbg}{RGB}{240, 255, 240}
\definecolor{solutionbg}{RGB}{242, 255, 242}
\definecolor{rulegray}{RGB}{180, 180, 180}
\definecolor{tableheadbg}{RGB}{0, 60, 113}

% ── Header/Footer ─────────────────────────────────────────────
\pagestyle{fancy}
\fancyhf{}
\renewcommand{\headrulewidth}{0pt}
\renewcommand{\footrulewidth}{0.4pt}
\renewcommand{\sectionmark}[1]{\markboth{#1}{}}
\fancyfoot[L]{\footnotesize\color{accent}\nouppercase{\leftmark}}
\fancyfoot[R]{\footnotesize\color{accent}\thepage\,/\,\pageref{LastPage}}
\fancypagestyle{plain}{%
  \fancyhf{}%
  \renewcommand{\headrulewidth}{0pt}%
  \renewcommand{\footrulewidth}{0.4pt}%
  \fancyfoot[L]{\footnotesize\color{accent}\nouppercase{\leftmark}}%
  \fancyfoot[R]{\footnotesize\color{accent}\thepage\,/\,\pageref{LastPage}}%
}

% ── Typografie ────────────────────────────────────────────────
\titleformat{\section}
    {\color{accent}\large\bfseries}{}{0pt}{}
    [\vspace{-4pt}\color{accent}\rule{\linewidth}{2pt}]
\titleformat{\subsection}
    {\normalsize\bfseries\color{accent!85!black}}{}{0pt}{}
    [\vspace{-3pt}\color{accent!40}\rule{\linewidth}{0.4pt}]
\titleformat{\subsubsection}
    {\small\bfseries\color{accent!70!black}}{}{0pt}{}{}
\titlespacing*{\section}{0pt}{12pt}{4pt}
\titlespacing*{\subsection}{0pt}{8pt}{2pt}
\titlespacing*{\subsubsection}{0pt}{6pt}{1pt}

\setlength{\parindent}{0pt}
\setlength{\parskip}{4pt}
\renewcommand{\arraystretch}{1.25}
\setlist[itemize]{noitemsep, topsep=3pt, parsep=0pt, leftmargin=1.4em}
\setlist[enumerate]{noitemsep, topsep=3pt, parsep=0pt, leftmargin=1.6em}
\setlist[itemize,2]{label=\textendash, leftmargin=1.4em}

% ── Boxen ──────────────────────────────────────────────────────
\newtcolorbox{figurebox}[1]{
  colback=figurebg, colframe=accent!50,
  title={Abbildung: #1}, fonttitle=\small\bfseries,
  breakable, before skip=6pt, after skip=6pt
}
\newtcolorbox{tablehintbox}[1]{
  colback=tablehintbg, colframe=green!50!black!50,
  title={Tabelle: #1}, fonttitle=\small\bfseries,
  breakable, before skip=6pt, after skip=6pt
}
\newtcolorbox{solutionbox}[1]{
  colback=solutionbg, colframe=green!60!black!60,
  title={#1}, fonttitle=\small\bfseries,
  breakable, before skip=6pt, after skip=6pt
}

\lstset{
  basicstyle=\ttfamily\small,
  backgroundcolor=\color{codebg},
  breaklines=true,
  frame=single, framerule=0pt,
  xleftmargin=4pt, xrightmargin=4pt,
  aboveskip=6pt, belowskip=6pt,
  keepspaces=true
}

% ── TOC-Look ──────────────────────────────────────────────────
\renewcommand{\contentsname}{\color{accent}Inhalt}

\providecommand{\nichttitle}{Lernzettel}
\providecommand{\nichtsubtitle}{}

\renewcommand{\nichttitle}{Relationen, Algebra, Diskrete Mathematik}
\renewcommand{\nichtsubtitle}{Lernzettel}


\begin{document}

\begin{center}
{\bfseries\Large\color{accent}\nichttitle}\\[2pt]
\color{accent}\rule{0.35\linewidth}{1pt}\color{black}
\end{center}
\vspace{2pt}

{\small\tableofcontents}
\newpage

% ── BODY ──────────────────────────────────────────────────────

\section{Logik}


\subsection{Grundlagen}

\begin{itemize}
  \item \textbf{Logik}: Wissen darstellen, Schlussfolgerungen ziehen. Anwendung: Hardware, Programmiersprachen, DBs, Algorithmen, ML.
  \item \textbf{Atomare Aussage}: Satz, eindeutig wahr (w) oder falsch (f). Verknüpfbar mit ∧, ∨, ¬.
  \item \textbf{Aussagenvariablen}: a, b, c, … (nicht w/f).
  \item \textbf{Belegung}: Zuordnung von Wahrheitswerten zu Variablen.
\end{itemize}


\subsubsection{Beispiele Aussage / keine Aussage}
\begin{itemize}
  \item "Berlin ist Hauptstadt Deutschlands" → wahr; "6+7=13" → wahr
  \item "Bonn wird wieder Hauptstadt" → keine (nicht eindeutig); "Na sowas!" → keine; "Dieser Satz ist falsch" → Paradoxon
\end{itemize}


\subsection{Operatoren}


{\small
\begin{tabularx}{\linewidth}{XXX}
\hline
\rowcolor{tableheadbg}
{\color{white}\textbf{Operator}} & {\color{white}\textbf{Symbol}} & {\color{white}\textbf{Bezeichnung}} \\[2pt]
\hline
\endhead
\hline
\endfoot
und & ∧ & Konjunktion \\
oder & ∨ & Disjunktion (nicht "entweder…oder") \\
nicht & ¬ & Negation \\
wenn-dann & → & Implikation \\
genau-dann-wenn & ↔ & Äquivalenz \\
\end{tabularx}
}


\begin{itemize}
  \item a ∧ b: wahr nur wenn beide wahr
  \item a ∨ b: wahr wenn mind. eine wahr
  \item a → b = ¬a ∨ b; wahr außer (w,f)
  \item a ↔ b = (a ∧ b) ∨ (¬a ∧ ¬b)
  \item Neutralelemente: w für ∧ (a ∧ w = a); f für ∨ (a ∨ f = a)
\end{itemize}


\subsection{Syntax / Semantik}

\textbf{Syntax} (Bildung korrekter Formeln):

\begin{enumerate}
  \item Atomare Aussagen xᵢ sowie w, f sind Formeln
  \item a, b Formeln → (a ∧ b), (a ∨ b) Formeln
  \item a Formel → ¬a Formel
\end{enumerate}

\begin{quote}
\textbf{Merke:} → und ↔ tauchen in Syntax nicht auf (durch ∧/∨/¬ darstellbar).
\end{quote}

\textbf{Semantik}: Abbildung aller Aussagevariablen einer Formel F(x₁,…,xₙ) auf \{w, f\}.


\begin{itemize}
  \item \textbf{Erfüllbar}: ∃ Belegung mit F = w
  \item \textbf{Tautologie}: alle Belegungen → F = w
  \item \textbf{Widerspruchsvoll}: keine Belegung → F = w
  \item F Tautologie ↔ ¬F widerspruchsvoll
  \item Bsp: x ∨ ¬x (Tautologie); x ∧ ¬x (widerspruchsvoll)
\end{itemize}


\subsection{Vorrangregeln}


{\small
\begin{tabularx}{\linewidth}{XX}
\hline
\rowcolor{tableheadbg}
{\color{white}\textbf{Operator}} & {\color{white}\textbf{Rang}} \\[2pt]
\hline
\endhead
\hline
\endfoot
( ) & 0 \\
¬ & 1 \\
∧ & 2 \\
∨ & 3 \\
→ & 4 \\
↔ & 5 \\
\end{tabularx}
}


\begin{itemize}
  \item Klammern/∧ weglassen wie Multiplikation: a ∧ b → ab
  \item Gleiche Verknüpfungen: links nach rechts
\end{itemize}


\subsection{Logische Gesetze}


{\small
\begin{tabularx}{\linewidth}{XXX}
\hline
\rowcolor{tableheadbg}
{\color{white}\textbf{Nr}} & {\color{white}\textbf{Name}} & {\color{white}\textbf{Tautologie}} \\[2pt]
\hline
\endhead
\hline
\endfoot
1 & Idempotenz & a ∧ a ↔ a; a ∨ a ↔ a \\
2 & Kommutativ & a ∧ b ↔ b ∧ a; a ∨ b ↔ b ∨ a \\
3 & Assoziativ & (a ∧ b) ∧ c ↔ a ∧ (b ∧ c); analog ∨ \\
4 & Absorption & a ∧ (a ∨ b) ↔ a; a ∨ (a ∧ b) ↔ a \\
5 & Distributiv & a ∧ (b ∨ c) ↔ (a ∧ b) ∨ (a ∧ c); a ∨ (b ∧ c) ↔ (a ∨ b) ∧ (a ∨ c) \\
6 & Doppelte Negation & ¬(¬a) ↔ a \\
7 & DeMorgan & ¬(a ∧ b) ↔ ¬a ∨ ¬b; ¬(a ∨ b) ↔ ¬a ∧ ¬b \\
\end{tabularx}
}



\subsection{Normalformen}

\begin{itemize}
  \item \textbf{Literal} xᵢ: Variable oder negierte Variable
  \item \textbf{n-stelliger Konjunktionsterm K}: x₁ ∧ … ∧ xₙ
  \item \textbf{n-stelliger Disjunktionsterm D}: x₁ ∨ … ∨ xₙ
  \item Duplikate entfernen; n=1 erlaubt
  \item \textbf{KN}: D₁ ∧ … ∧ Dₘ; \textbf{DN}: K₁ ∨ … ∨ Kₘ
  \item \textbf{kDN/kKN} (kanonisch): jeder Term enthält ALLE Variablen → Minterme/Maxterme; eindeutig
\end{itemize}


\subsubsection{Verfahren DN/KN}
\begin{enumerate}
  \item ↔, → ersetzen durch ¬, ∨, ∧
  \item DeMorgan, bis Negationen nur vor Variablen
  \item Distributivgesetz (DN: ∧ über ∨; KN: ∨ über ∧)
  \item Idempotenz: gleiche Terme zusammenfassen
  \item Neutralelemente: ¬a∧a=f; ¬a∨a=w; a∧w=a; a∨f=a
\end{enumerate}


\subsubsection{Verfahren kDN/kKN}
\begin{enumerate}
  \item Term in DN/KN suchen, in dem Variable a fehlt
  \item K → K ∧ (a ∨ ¬a) bzw. D → D ∨ (a ∧ ¬a)
  \item Distributiv anwenden, Idempotenz
  \item Bsp: f(a,b) = a → a ∧ (b ∨ ¬b) → ab ∨ a¬b
\end{enumerate}


\subsubsection{Sortierung}
\begin{itemize}
  \item Innerhalb K/D: lexikographisch / nach Index aufsteigend
  \item Zwischen Termen: ¬x:=0, x:=1 → Binärzahl, aufsteigend
\end{itemize}


\subsubsection{Äquivalenz}
Umformung in Normalform — gleiche Normalform = äquivalent.



\subsection{Logisches Schließen}

\begin{itemize}
  \item \{(a→b),(b→c)\} ⊨ (a→c) — Voraussetzungen liefern Schlussfolgerung
  \item Grundlage: (a→b) ∧ (b→c) → (a→c) ist Tautologie
  \item \textbf{Kontraposition}: a→b ↔ ¬b→¬a; a→b ≠ ¬a→¬b
\end{itemize}


\subsection{Resolution}

\begin{itemize}
  \item \textbf{Beweis durch Widerspruch}: zeige (A₁ ∧ … ∧ Aₙ ∧ ¬b) widerspruchsvoll
  \item Voraussetzung: KN
  \item \textbf{Klausel}: Disjunktion von Literalen, geschrieben als Menge
  \item KN → Klauselmenge: ∧ und ∨ entfernen, Mengen-Notation
  \item \textbf{Resolvente}: A enthält a, B enthält ¬a → res\{A,B\} = (A∖\{a\}) ∪ (B∖\{¬a\})
  \item Komplementäre Literale in Resolvente → verwerfen
  \item \textbf{Leere Klausel} □: Widerspruch bewiesen
  \item Pro Schritt nur EIN Literal entfernen
  \item Einzelne Klausel kann nicht resolviert werden
  \item Verwendete Klauseln bleiben in Menge
  \item Kein weiterer Resolutionsschritt möglich und kein □ → Behauptung falsch
\end{itemize}


\subsubsection{Beispiel: \{(a→b),(b→c)\} ⊨ (a→c)}
\begin{itemize}
  \item KN: (¬a∨b) ∧ (¬b∨c) ∧ a ∧ ¬c
  \item Klauselmenge: \{\{¬a,b\},\{¬b,c\},\{a\},\{¬c\}\}
  \item res(\{¬a,b\},\{¬b,c\})=\{¬a,c\} → res(.,a)=\{c\} → res(c,¬c)=□
\end{itemize}

\begin{quote}
\textbf{Merke:} Resolution für ↔-Vermutungen unhandlich → SAT-Algorithmen.
\end{quote}


\subsection{Prädikatenlogik}

\begin{itemize}
  \item \textbf{Individuen}: Objekte mit Eigenschaften
  \item \textbf{n-stelliges Prädikat} P(x₁,…,xₙ): beschreibt Eigenschaft
  \item \textbf{Quantoren}: ∀ (Allquantor), ∃ (Existenzquantor)
  \item \textbf{Funktionsterme}: liefern Individuen
\end{itemize}


\subsubsection{Beispiele}
\begin{itemize}
  \item ∀x ∈ ℕ (Quadrat(x) ≥ x)
  \item "Alle Menschen sind sterblich": ∀x(Mensch(x) → sterblich(x))
  \item "Sokrates ist Mensch": Mensch(Sokrates) = w
  \item "Es gibt keinen Menschen, der >20m springt": ¬∃x(Mensch(x) ∧ springt\_weiter(x,20))
\end{itemize}

Resolution übertragbar (Substitution, Unifikation, Umbenennungen).



\section{Mengen}


\subsection{Begriff \& Schreibweisen}

\textbf{Menge (Cantor)}: Zusammenfassung unterscheidbarer Objekte; Objekte = Elemente. Mengen Großbuchstaben (M), Elemente klein (x).


\begin{itemize}
  \item Aufzählend: M = \{1, 2, 3\}
  \item Beschreibend: M = \{x | x ∈ ℕ ∧ 1 ≤ x ≤ 3\}
  \item Charakteristische Aussageform: M = \{x | A(x)\}
  \item Grafisch: Venn-Diagramm
  \item Ungültig: \{1,2,3,3\} (nicht unterscheidbar), Verwendung anderer Klammern als \{\}
\end{itemize}


\subsection{Zugehörigkeit \& Mächtigkeit}

\begin{itemize}
  \item a ∈ M / a ∉ M
  \item |M|: Anzahl Elemente
  \item |∅| = 0, |\{∅\}| = 1, |P(∅)| = 1, |P(\{∅\})| = 2
\end{itemize}


\subsection{Leere Menge ∅}

\begin{itemize}
  \item Schreibweise: ∅, \{\}
  \item Teilmenge jeder Menge
  \item Nicht Element jeder Menge (nur Element jeder Potenzmenge)
\end{itemize}


\subsection{Gleichheit \& Teilmenge}

\begin{itemize}
  \item Extensionalität: M = N ↔ ∀x(x ∈ M ↔ x ∈ N); Reihenfolge irrelevant
  \item Echte Teilmenge: T ⊂ M (T ≠ M)
  \item Unechte Teilmenge: T ⊆ M
  \item Obermenge: M ⊇ T
\end{itemize}


\subsection{Zahlenmengen \& Intervalle}


{\small
\begin{tabularx}{\linewidth}{XX}
\hline
\rowcolor{tableheadbg}
{\color{white}\textbf{Menge}} & {\color{white}\textbf{Definition}} \\[2pt]
\hline
\endhead
\hline
\endfoot
ℕ & \{1,2,3,…\} \\
ℕ₀ & \{0,1,2,…\} \\
ℤ & \{…,−1,0,1,…\} \\
ℚ & \{p/q \\
ℝ & alle Dezimalbrüche a₀.a₁a₂…, a₀∈ℤ, aᵢ∈\{0,…,9\} \\
ℂ & komplexe Zahlen \\
\end{tabularx}
}


\begin{itemize}
  \item [a,b]: abgeschlossen; (a,b): offen; (a,b], [a,b): halboffen
\end{itemize}


\subsection{Potenzmenge \& Komplement}

\begin{itemize}
  \item \textbf{P(M)} = \{x | x ⊆ M\}; |P(M)| = 2\textasciicircum{}|M|
  \item \textbf{Grundmenge E}: enthält alle betrachteten Mengen
  \item \textbf{Komplement}: C\_E(M) = M̄ = E ∖ M (M ⊆ E)
  \item Bsp: E=\{1,…,5\}, M=\{2,4\} → C\_E(M)=\{1,3,5\}
\end{itemize}


\subsection{Mengenalgebra}


{\small
\begin{tabularx}{\linewidth}{XX}
\hline
\rowcolor{tableheadbg}
{\color{white}\textbf{Operation}} & {\color{white}\textbf{Definition}} \\[2pt]
\hline
\endhead
\hline
\endfoot
Schnitt & M ∩ N = \{x \\
Vereinigung & M ∪ N = \{x \\
Differenz & M ∖ N = \{x \\
Disjunkt & A ∩ B = ∅ \\
Kartesisches Produkt & M × N = \{(x,y) \\
\end{tabularx}
}



\subsubsection{Gesetze (analog Logik)}


{\small
\begin{tabularx}{\linewidth}{XX}
\hline
\rowcolor{tableheadbg}
{\color{white}\textbf{Name}} & {\color{white}\textbf{Gesetz}} \\[2pt]
\hline
\endhead
\hline
\endfoot
Idempotenz & M ∩ M = M; M ∪ M = M \\
Kommutativ & M ∩ N = N ∩ M; analog ∪ \\
Assoziativ & (M∩N)∩W = M∩(N∩W) \\
Absorption & M ∩ (M∪N) = M; M ∪ (M∩N) = M \\
Distributiv & M ∩ (N∪W) = (M∩N) ∪ (M∩W) \\
Doppeltes Komplement & C\_E(C\_E(M)) = M \\
DeMorgan & C\_E(M∩N) = C\_E(M) ∪ C\_E(N); C\_E(M∪N) = C\_E(M) ∩ C\_E(N) \\
\end{tabularx}
}



\section{Algebraische Strukturen}

\textbf{Algebraische Struktur}: Menge M + Verknüpfungen + Axiome. Verknüpfung v: M × M → M.



\subsection{Gruppe (G, +)}
\begin{enumerate}
  \item Abgeschlossenheit: a+b ∈ G
  \item Assoziativ: (a+b)+c = a+(b+c)
  \item Neutrales Element e: a+e = a
  \item Inverses a\textit{: a+a} = e
\end{enumerate}

\textbf{Kommutative Gruppe}: zusätzlich a+b = b+a.

\begin{itemize}
  \item (ℝ,+), (ℤ,+) kommutative Gruppen
  \item (ℕ₀,+) keine Gruppe (kein Inverses)
  \item (ℝ∖\{0\}, ∗) kommutative Gruppe
\end{itemize}


\subsection{Körper (K, +, ∗)}
\begin{enumerate}
  \item + und ∗ definiert
  \item (K, +) kommutative Gruppe
  \item (K∖\{0\}, ∗) kommutative Gruppe
  \item Distributiv: a∗(b+c) = (a∗b)+(a∗c)
\end{enumerate}

\begin{itemize}
  \item (ℝ,+,∗) ist Körper
\end{itemize}


\section{Relationen}


\subsection{Begriff}

\textbf{Relation}: Beziehung zwischen Elementen — besteht / besteht nicht.


\begin{itemize}
  \item 2-stellig: R ⊆ M₁ × M₂; bei M=M₁=M₂: R ⊆ M²
  \item Schreibweisen: (a,b) ∈ R, aRb, R(a,b)
  \item Mathematische Operatoren <, >, ≤, ≥, =, ≠, ≈ sind Relationen
  \item n-stellig: R ⊆ M₁ × … × Mₙ — Tupelmenge
  \item Datenbankrelation = n-stellige Relation; Tabellenkopf = Schema
\end{itemize}


\subsection{Eigenschaften}


{\small
\begin{tabularx}{\linewidth}{XX}
\hline
\rowcolor{tableheadbg}
{\color{white}\textbf{Eigenschaft}} & {\color{white}\textbf{Bedingung}} \\[2pt]
\hline
\endhead
\hline
\endfoot
Reflexiv & ∀x: xRx \\
Irreflexiv & ∄x: xRx \\
Symmetrisch & xRy → yRx \\
Asymmetrisch & xRy → ¬yRx \\
Antisymmetrisch & xRy ∧ yRx → x=y \\
Transitiv & xRy ∧ yRz → xRz \\
Linear & ∀x,y: xRy ∨ yRx \\
Konnex & x≠y → xRy ∨ yRx \\
\end{tabularx}
}


\begin{quote}
\textbf{Merke:} asymmetrisch ⇒ antisymmetrisch.
\end{quote}


\subsection{Äquivalenz- \& Ordnungsrelationen}

\textbf{Äquivalenzrelation}: reflexiv + symmetrisch + transitiv. Teilt Grundmenge in Äquivalenzklassen.



{\small
\begin{tabularx}{\linewidth}{XXX}
\hline
\rowcolor{tableheadbg}
{\color{white}\textbf{Ordnung}} & {\color{white}\textbf{Irreflexiv}} & {\color{white}\textbf{Reflexiv}} \\[2pt]
\hline
\endhead
\hline
\endfoot
Quasiordnung & irreflex. + transitiv & reflex. + transitiv \\
Halbordnung & + asymmetrisch & + antisymmetrisch (Bsp. ≤) \\
Vollordnung & + konnex (Bsp. <) & + linear (Bsp. ≤) \\
\end{tabularx}
}


\begin{quote}
\textbf{Merke:} "größer oder gleich groß" ist keine reflex. Halbordnung (nicht antisymmetrisch).
\end{quote}


\section{Abbildungen / Funktionen}

\textbf{Abbildung = Funktion}: f: D → W, x ↦ f(x). Spezielle 2-stellige Relation.

\begin{itemize}
  \item D: Definitionsmenge, W: Wertemenge; y: Bild, x: Urbild
  \item \textbf{Linksvollständig}: jedem x ∈ D ist etwas zugeordnet
  \item \textbf{Rechtseindeutig}: höchstens ein Bild pro x
\end{itemize}


{\small
\begin{tabularx}{\linewidth}{XXX}
\hline
\rowcolor{tableheadbg}
{\color{white}\textbf{Eigenschaft}} & {\color{white}\textbf{Definition}} & {\color{white}\textbf{Anschauung}} \\[2pt]
\hline
\endhead
\hline
\endfoot
Injektiv & f(x₁)=f(x₂) → x₁=x₂ & linkseindeutig; max. 1 Pfeil pro y \\
Surjektiv & ∀y∈W ∃x: f(x)=y & rechtsvollständig; ≥1 Pfeil pro y \\
Bijektiv & injektiv ∧ surjektiv & invertierbar; genau 1 Pfeil pro y \\
\end{tabularx}
}


\begin{itemize}
  \item Bsp: f:ℝ→ℝ, f(x)=x² → weder injektiv noch surjektiv
  \item f:ℝ⁺→ℝ⁺, f(x)=x² → bijektiv
  \item Bewertung hängt von gewählten D, W ab
\end{itemize}


\section{Relationale Datenbanken}

\begin{itemize}
  \item \textbf{Relation}: Tabelle
  \item \textbf{Tupel}: Datensatz
  \item \textbf{Schema}: Kopf
  \item \textbf{Primärschlüssel}: Attribut(kombination), eindeutig + minimal
  \item \textbf{Fremdschlüssel}: in anderer Relation Primärschlüssel; stellt Beziehungen her
\end{itemize}


\subsection{Probleme schlechter Entwürfe}

\begin{itemize}
  \item Alles in einer Tabelle → Redundanz, Update-Anomalien
  \item Spalte je Sportart → viele 0-Felder, neue Sportart = neue Spalte
  \item \textbf{Lösung}: Normalisierung (zwei Tabellen, FK-Verknüpfung)
\end{itemize}


\subsection{Relationenalgebra}

\begin{itemize}
  \item \textbf{Mengenoperationen} ∩, ∪, ∖: nur bei kompatiblem Schema; ∪ ohne Duplikate
  \item \textbf{Kreuzprodukt} ×: ohne Kompatibilität, jedes Tupel mit jedem
  \item \textbf{Selektion} σ\_F(r): Zeilen, die F erfüllen — SQL: \texttt{where}
  \item \textbf{Projektion} Π\_Y(r): Spalten Y; mengentheoretisch ohne Duplikate — SQL: \texttt{select distinct}
  \item \textbf{Join} ⋈: Verknüpfung über gleiche Spalten via Primär-/Fremdschlüssel
  \item \textbf{Natural Join}: Spalten gleichen Namens \& Inhalts
  \item \textbf{SQL inner join}: \texttt{from A,B where A.k=B.k}
  \item \textbf{Weitere Joins}: left, right, outer, equi, semi
\end{itemize}


\section{Boolesche Algebra}

\textbf{Boolesche Algebra} (M; ⊗, ⊕, ∼): Grundlage digitaler Schaltungen.

\begin{itemize}
  \item Menge M ≥ 2 Elemente
  \item ⊗: M×M→M (Produkt), ⊕: M×M→M (Summe), ∼: M→M (Komplement)
\end{itemize}


\subsection{Axiome}


{\small
\begin{tabularx}{\linewidth}{XXX}
\hline
\rowcolor{tableheadbg}
{\color{white}\textbf{Gesetz}} & {\color{white}\textbf{Form 1}} & {\color{white}\textbf{Form 2}} \\[2pt]
\hline
\endhead
\hline
\endfoot
Kommutativ & a⊗b = b⊗a & a⊕b = b⊕a \\
Distributiv & a⊗(b⊕c) = (a⊗b)⊕(a⊗c) & a⊕(b⊗c) = (a⊕b)⊗(a⊕c) \\
Neutralelemente & ∃1: a⊗1 = a & ∃0: a⊕0 = a \\
Komplement & a⊗∼a = 0 & a⊕∼a = 1 \\
\end{tabularx}
}


\textbf{Dualitätsprinzip}: ⊗↔⊕, 0↔1; jeder Beweis gilt dual.



\subsection{Abgeleitete Gesetze}


{\small
\begin{tabularx}{\linewidth}{XXX}
\hline
\rowcolor{tableheadbg}
{\color{white}\textbf{Name}} & {\color{white}\textbf{Form 1}} & {\color{white}\textbf{Form 2}} \\[2pt]
\hline
\endhead
\hline
\endfoot
Idempotenz & a⊗a = a & a⊕a = a \\
Assoziativ & (a⊗b)⊗c = a⊗(b⊗c) & (a⊕b)⊕c = a⊕(b⊕c) \\
Absorption & a⊗(a⊕b) = a & a⊕(a⊗b) = a \\
Doppeltes Komplement & ∼(∼a) = a & — \\
DeMorgan & ∼(a⊗b) = ∼a ⊕ ∼b & ∼(a⊕b) = ∼a ⊗ ∼b \\
Neutrale komplementär & ∼0 = 1 & ∼1 = 0 \\
0-1-Gesetze & a⊗0 = 0 & a⊕1 = 1 \\
\end{tabularx}
}



\subsubsection{Beweis a⊗0 = 0}
\begin{enumerate}
  \item a⊗∼a = 0
  \item a⊗(∼a ⊕ 0) = 0
  \item (a⊗∼a) ⊕ (a⊗0) = 0
  \item 0 ⊕ (a⊗0) = 0
  \item a⊗0 = 0
\end{enumerate}


\subsection{Modelle}

\begin{itemize}
  \item \textbf{Mengenalgebra}: P(M) als Träger; ⊗≙∩, ⊕≙∪, ∼≙Komplement; Neutrale: M, ∅
  \item \textbf{Schaltalgebra}: M = B = \{0,1\}; zweielementig; Schalter offen=0/zu=1
\end{itemize}


\subsection{Schaltalgebra-Operatoren}


{\small
\begin{tabularx}{\linewidth}{XXX}
\hline
\rowcolor{tableheadbg}
{\color{white}\textbf{Operator}} & {\color{white}\textbf{Realisierung}} & {\color{white}\textbf{Lampe brennt wenn}} \\[2pt]
\hline
\endhead
\hline
\endfoot
a ∧ b & Reihenschaltung & beide zu \\
a ∨ b & Parallelschaltung & mind. 1 zu \\
a′ & Ruhekontakt & offen \\
\end{tabularx}
}



{\small
\begin{tabularx}{\linewidth}{XXXX}
\hline
\rowcolor{tableheadbg}
{\color{white}\textbf{a}} & {\color{white}\textbf{b}} & {\color{white}\textbf{a∧b}} & {\color{white}\textbf{a∨b}} \\[2pt]
\hline
\endhead
\hline
\endfoot
0 & 0 & 0 & 0 \\
0 & 1 & 0 & 1 \\
1 & 0 & 0 & 1 \\
1 & 1 & 1 & 1 \\
\end{tabularx}
}



{\small
\begin{tabularx}{\linewidth}{XX}
\hline
\rowcolor{tableheadbg}
{\color{white}\textbf{a}} & {\color{white}\textbf{a′}} \\[2pt]
\hline
\endhead
\hline
\endfoot
0 & 1 \\
1 & 0 \\
\end{tabularx}
}



\subsection{Boolesche Terme}

Syntax:

\begin{enumerate}
  \item Konstanten/Variablen sind Terme
  \item A, B Terme → A′, (A∧B), (A∨B)
  \item nur endlich viele Schritte
\end{enumerate}

Vorrang: ( ), ′, ∧, ∨ (1–3); ∧ weglassbar.



\subsection{Boolesche Funktion}

\textbf{f: Bⁿ → B}: 2ⁿ Belegungen → 2\textasciicircum{}(2ⁿ) mögliche Funktionen bei n Variablen.



\subsubsection{Normalformen aus Wertetabelle}
\begin{itemize}
  \item \textbf{kDN}: alle Zeilen mit f=1; Minterm pro Zeile (a=0→a′, a=1→a); Minterme ∨-verbinden
\begin{itemize}
  \item Zeile 010 → a′bc′
\end{itemize}
  \item \textbf{kKN}: alle Zeilen mit f=0; Maxterm (a=0→a, a=1→a′); ∧-verbinden
\begin{itemize}
  \item Zeile 010 → a ∨ b′ ∨ c
\end{itemize}
\end{itemize}


\subsection{Minimierung — KV-Diagramme (Karnaugh-Veitch)}

\begin{itemize}
  \item Spalten/Zeilen: Variablenbelegungen im Gray-Code (00, 01, 11, 10) — 1 Bit Änderung pro Schritt
  \item Zellen = Funktionswerte; nur Einsen eintragen
  \item Einserblöcke maximal groß identifizieren; dürfen über Kanten hinausgehen
  \item Variable in Block durchgängig konstant → bleibt; wechselt → eliminiert
  \item \textbf{Don't-care (d)}: irrelevant — beliebig nutzbar für größere Blöcke
  \item Beispiel 4-Var: Block über Kante b′d′ → DM: f = b′d′
  \item Beispiel 4-Var: Zeilen c=0,d=1 + a=1,d=1 → DM: f = ad ∨ c′d
\end{itemize}

\begin{quote}
\textbf{Merke:} Skaliert nicht: 2³² ≈ 4,3·10⁹ Zeilen.
\end{quote}

Weitere Methode: \textbf{Quine-McCluskey} (komplexere Fälle).



\section{Schaltnetze}

\textbf{Schaltnetz}: F: Bⁿ → Bᵐ; verarbeitet n Eingänge → m Ausgänge.



\subsection{Halbaddierer (HA)}


{\small
\begin{tabularx}{\linewidth}{XXXX}
\hline
\rowcolor{tableheadbg}
{\color{white}\textbf{a}} & {\color{white}\textbf{b}} & {\color{white}\textbf{s}} & {\color{white}\textbf{ü}} \\[2pt]
\hline
\endhead
\hline
\endfoot
0 & 0 & 0 & 0 \\
0 & 1 & 1 & 0 \\
1 & 0 & 1 & 0 \\
1 & 1 & 0 & 1 \\
\end{tabularx}
}


\begin{itemize}
  \item s = a′b ∨ ab′ = a ⊕ b
  \item ü = ab
\end{itemize}


\subsection{Volladdierer (VA)}


{\small
\begin{tabularx}{\linewidth}{XXXXX}
\hline
\rowcolor{tableheadbg}
{\color{white}\textbf{a}} & {\color{white}\textbf{b}} & {\color{white}\textbf{ü₁}} & {\color{white}\textbf{s}} & {\color{white}\textbf{ü₂}} \\[2pt]
\hline
\endhead
\hline
\endfoot
0 & 0 & 0 & 0 & 0 \\
0 & 0 & 1 & 1 & 0 \\
0 & 1 & 0 & 1 & 0 \\
0 & 1 & 1 & 0 & 1 \\
1 & 0 & 0 & 1 & 0 \\
1 & 0 & 1 & 0 & 1 \\
1 & 1 & 0 & 0 & 1 \\
1 & 1 & 1 & 1 & 1 \\
\end{tabularx}
}


\begin{itemize}
  \item s = a ⊕ b ⊕ ü₁
  \item ü₂ = ab ∨ aü₁ ∨ bü₁
\end{itemize}


\subsection{Anwendung Spülmaschinensteuerung}

f(a,b,c) = 1, wenn Tür offen ODER (Wasser zu niedrig UND Heizung an).



\section{Graphen}


\subsection{Grundlagen}

\textbf{Graph}: G = (V, E) bzw. G = (V, E, W); V = Knoten, E = Kanten, W = Gewichte.


\begin{itemize}
  \item \textbf{Knoten}: Objekte (Städte, Personen, Computer)
  \item \textbf{Kanten}: Beziehungen; können gerichtet sein
  \item \textbf{Gewicht}: Kosten/Distanz/Wahrscheinlichkeit; bei gerichteten Graphen richtungsabhängig
  \item \textbf{Knotengrad}: Anzahl angrenzender Kanten
  \item Knotenposition irrelevant (außer bei räumlicher Verortung)
\end{itemize}


\subsection{Typen}


{\small
\begin{tabularx}{\linewidth}{XX}
\hline
\rowcolor{tableheadbg}
{\color{white}\textbf{Typ}} & {\color{white}\textbf{Merkmal}} \\[2pt]
\hline
\endhead
\hline
\endfoot
Ungerichtet / gerichtet & mit/ohne Kantenrichtung \\
Mit Gewichten & Kanten tragen Werte \\
Multigraph & Mehrfachkanten zwischen 2 Knoten \\
Leer & keine Kanten \\
Vollständig & Kanten zwischen allen Knoten \\
Baum & keine Zyklen \\
Kreis & \textbackslash{} \\
Stern & 1 Zentrum mit allen verbunden \\
Bipartit & 2 Mengen, Kanten nur zwischen ihnen \\
\end{tabularx}
}



\subsection{Darstellungsformen}

\begin{itemize}
  \item \textbf{Zeichnung}
  \item \textbf{Adjazenzmatrix} A (n×n): aᵢⱼ ≠ 0 wenn Kante; Gewicht oder 0/1; Multigraph ohne Gewicht: aᵢⱼ = \#Kanten; ungerichtet → symmetrisch
  \item \textbf{Adjazenzliste}: pro Knoten Liste benachbarter Knoten; gerichtet: 2 Listen (Vor-/Nachfolger); effizient bei spärlichen Graphen
\end{itemize}


\subsection{Pfade und Zyklen}

\begin{itemize}
  \item \textbf{Pfad}: Sequenz von Kanten; jeder Knoten höchstens einmal
  \item \textbf{Zyklus}: geschlossener Pfad mit unterschiedlichen Kanten; Bäume zyklenfrei
  \item \textbf{Eulerscher Pfad/Zyklus}: jede Kante genau einmal
  \item \textbf{Hamiltonscher Pfad/Zyklus}: jeder Knoten genau einmal
\end{itemize}


\section{Kürzeste Pfade}

\textbf{Definition}: Pfad zwischen zwei Knoten mit minimaler Gewichtssumme.



\subsection{Lösungsstrategien}

\begin{itemize}
  \item \textbf{Brute Force}: alle Touren prüfen — optimal, aber exponentieller Aufwand (Bsp: 7 Knoten/12 Kanten → 29 Pfade)
  \item \textbf{Greedy / Best-first}: nächste zulässige Kante mit geringstem Gewicht — oft suboptimal, Sackgassen möglich
\end{itemize}


\subsection{Dijkstra-Algorithmus}

\textbf{Zweck}: Kürzeste Pfade von Startknoten s zu allen anderen Knoten.


\textbf{Variablen}: d (Distanz), π (Status), α/previous (Vorgänger), u (aktuell), w(u,v) (Gewicht).


\textbf{Einschränkungen}:

\begin{itemize}
  \item Funktioniert NICHT mit negativen Gewichten
  \item Kein Zielknoten als Eingabe → ineffizient für Punkt-zu-Punkt (Abbruch möglich, sobald Zielknoten visited)
\end{itemize}


\subsubsection{Initialisierung}
\begin{lstlisting}
for v ∈ V: v.d = ∞; v.π = unvisited; v.α = none
s.d = 0
\end{lstlisting}



\subsubsection{Update}
\begin{lstlisting}
update(u,v,w):
  if v.d > u.d + w(u,v):
    v.d = u.d + w(u,v); v.α = u
\end{lstlisting}



\subsubsection{Hauptschleife}
\begin{lstlisting}
u = s
while ∃ v: v.π == unvisited (oder Zielknoten):
  u.π = visited
  for jeden unvisited Nachbar v: update(u,v,w)
  u = unvisited Knoten mit minimalem d
\end{lstlisting}



\subsubsection{Sprachlich}
\begin{enumerate}
  \item d(s)=0, alle anderen ∞; alle unvisited
  \item Aktueller = s
  \item Solange nicht alle besucht: aktuellen markieren; für unbesuchte Nachbarn Distanz über aktuellen prüfen und ggf. updaten; nächster aktueller = unbesuchter mit kleinster d
\end{enumerate}


\subsection{A*-Algorithmus}

\textbf{Zweck}: Effiziente Pfadsuche s→z mittels Heuristik.


\textbf{Unterschied zu Dijkstra}:

\begin{itemize}
  \item Zusätzlich Priorität p = d + h(v,z)
  \item h(v,z): Heuristik (z. B. Luftlinie zu z)
  \item Knotenwahl nach minimalem p (nicht d)
\end{itemize}

\textbf{Nachteil}: falsche/überschätzende Heuristik → suboptimale Lösung.



\subsubsection{Initialisierung}
\begin{lstlisting}
for v ∈ V: v.d=∞; v.p=∞; v.π=unvisited; v.α=none
s.d=0; s.p=h(s,z)
\end{lstlisting}



\subsubsection{Update}
\begin{lstlisting}
update(u,v,w,z,h):
  if v.d > u.d + w(u,v):
    v.d = u.d + w(u,v); v.α = u; v.p = v.d + h(v,z)
\end{lstlisting}



\subsubsection{Hauptschleife}
\begin{lstlisting}
u=s
while z.π == unvisited:
  u.π = visited
  for unvisited Nachbarn v: update(u,v,w,z,h)
  u = unvisited Knoten mit minimalem p
\end{lstlisting}



\subsection{Vergleich}


{\small
\begin{tabularx}{\linewidth}{XX}
\hline
\rowcolor{tableheadbg}
{\color{white}\textbf{Algorithmus}} & {\color{white}\textbf{Eigenschaft}} \\[2pt]
\hline
\endhead
\hline
\endfoot
Dijkstra & 1 Startknoten; gerichtet \& ungerichtet; keine negativen Gewichte \\
Bellman–Ford & 1 Startknoten; negative Gewichte erlaubt \\
A* & Start + Ziel; Heuristik \\
Floyd–Warshall & Alle Paare; gerichtet, dichte Graphen \\
Johnson & Alle Paare; gerichtet, spärliche Graphen \\
\end{tabularx}
}



\subsection{R-Implementation (Datenstrukturen, gemeinsam für Dijkstra \& A*)}

\begin{itemize}
  \item \texttt{graph[x,y]} = Gewicht oder 0
  \item \texttt{distances} (Inf), \texttt{visited} (FALSE), \texttt{previous} (NA), \texttt{priorities} (Inf, A*), \texttt{heuristic}
\end{itemize}


\subsubsection{\texttt{dijkstra(graph, start, end=FALSE)}}
\begin{lstlisting}
distances=rep(Inf,n); visited=rep(FALSE,n); previous=rep(NA,n)
distances[start]=0; previous[start]=0; current=start
while(any(!visited)):
  visited[current]=TRUE
  if(current==end) break
  for i Nachbarn:
    if graph[current,i]!=0 && !visited[i] && distances[i] > distances[current]+graph[current,i]:
      distances[i]=distances[current]+graph[current,i]
      previous[i]=current
  current = unvisited Knoten mit kleinstem distances; sonst current=-1
\end{lstlisting}

\begin{itemize}
  \item \texttt{end=FALSE}: kürzeste Pfade zu allen Knoten
\end{itemize}


\subsubsection{\texttt{A*(graph, start, end, heuristic)}}
\begin{itemize}
  \item Wie Dijkstra, plus \texttt{priorities[start]=heuristic[start]}
  \item Update setzt zusätzlich \texttt{priorities[i] = distances[i] + heuristic[i]}
  \item Knotenwahl: \texttt{min priorities[i]} über \texttt{!visited[i]}
\end{itemize}


\subsubsection{\texttt{getpath(previous, start, end)} — Pfadrekonstruktion}
\begin{lstlisting}
i=end; path=end
while(i != 0):
  ii=previous[i]; i=ii
  if(i!=0) path=c(ii,path)
return path
\end{lstlisting}

\begin{itemize}
  \item Termination: \texttt{previous[start]=0}
\end{itemize}


\subsubsection{Wrapper \texttt{allpaths(graph)} — alle Paare}
\begin{lstlisting}
for s in 1:n:
  results[[s]] = dijkstra(graph, start=s, end=FALSE)
\end{lstlisting}

\begin{itemize}
  \item n Aufrufe (statt n²); ein Lauf liefert kürzeste Pfade von s zu allen anderen Knoten
\end{itemize}


\subsubsection{Vorlesungsbeispiel Ravensburg–Immenstaad (6 Knoten, ungerichtet)}


{\footnotesize
\begin{tabularx}{\linewidth}{XXXXXXX}
\hline
\rowcolor{tableheadbg}
{\color{white}\textbf{}} & {\color{white}\textbf{1}} & {\color{white}\textbf{2}} & {\color{white}\textbf{3}} & {\color{white}\textbf{4}} & {\color{white}\textbf{5}} & {\color{white}\textbf{6}} \\[2pt]
\hline
\endhead
\hline
\endfoot
1 & 0 & 1 & 0 & 0 & 4 & 0 \\
2 & 1 & 0 & 1 & 0 & 1 & 0 \\
3 & 0 & 1 & 0 & 1 & 0 & 0 \\
4 & 0 & 0 & 1 & 0 & 6 & 1 \\
5 & 4 & 1 & 0 & 6 & 0 & 0 \\
6 & 0 & 0 & 0 & 1 & 0 & 0 \\
\end{tabularx}
}


\begin{itemize}
  \item start=1, end=6
  \item A*-Heuristik: c(1,0,0,0,1000,0) → Knoten 5 wird durch hohen Wert gemieden
\end{itemize}


\subsubsection{Visualisierung (R)}
\begin{itemize}
  \item Pakete: \texttt{network}, \texttt{ggnet}, \texttt{GGally}
  \item \texttt{network(m, directed=F, names.eval="weights", ignore.eval=FALSE)}
  \item \texttt{ggnet2(...,mode="fruchtermanreingold",size=10,label=T,edge.label="weights",arrow.gap=0.05)}
  \item Pfadknoten rot; \texttt{set.seed(43)} für Reproduzierbarkeit
\end{itemize}


\section{TSP – Travelling Salesperson Problem}

\textbf{Definition}: n Städte alle genau einmal besuchen, zur Startstadt zurück, minimale Tourlänge.

\begin{itemize}
  \item Graph: V=Städte, E=Strecken, W=Längen
\end{itemize}


\subsection{Anwendungen}
\begin{itemize}
  \item Logistik (VRP-Erweiterungen)
  \item Mikrochip-Design, Leiterplattenbohrung
  \item Genom-Sequenzierung (Knoten = DNA-Teilstränge)
\end{itemize}


\subsection{Verwandte Probleme}
\begin{itemize}
  \item TSP ohne Rückkehr
  \item VRP (mehrere Reisende)
  \item Traveling Thief (TSP + Knapsack)
  \item Stochastisch / dynamisch
\end{itemize}


\subsection{Lösungsverfahren}

\textbf{Brute Force}: alle Touren durchprobieren — optimal, exponentielle Laufzeit.


\textbf{Nearest Neighbour Heuristik (NNA)}: Greedy-Verfahren ausgehend vom Startknoten.

\begin{enumerate}
  \item Start bei Startknoten
  \item Nächster Knoten: günstigste Kante zu unbesuchtem Knoten
  \item Wiederhole bis alle besucht
  \item Startknotenabhängig, meist suboptimal
  \item Verbesserung: jeden Knoten als Start probieren → bestes Ergebnis (Laufzeit ×n, weiterhin keine Garantie)
\end{enumerate}

\textbf{2-Opt Heuristik}: Lokale Suche zur Verbesserung einer vorhandenen Tour.

\begin{enumerate}
  \item Approximative Startlösung (z. B. NNA)
  \item Iterativ Teilstück invertieren, falls neue Tour kürzer → übernehmen
  \item Nächster Tausch
  \item Approximativ
\end{enumerate}

\textbf{Weitere Verfahren} (nicht behandelt): lineares Programmieren, evolutionäre Algorithmen.





% ──────────────────────────────────────────────────────────────

\end{document}
